\documentclass[authordate, editorial]{jote-new-article}

\usepackage{caption}

\usepackage{tabularx}

\usepackage{graphicx}

\usepackage{hyperref}

\usepackage[backend=biber,style=apa]{biblatex}


\jotetitle{Novel Cell Motility of Protozoa Living in the Hindgut of Termites}
\keywordsabstract{protozoan, rotary motor, fluid membranes, bacterial flagella}
\abstracttext{A symbiotic protozoan in the hindgut of termites was accidentally discovered. Its head and body continuously rotated in opposite directions without stopping or reversing directions. Rotations included the plasma membrane of each part of the cell, and the membrane was continuous across the surface shear zone. This protozoan offers direct visual evidence for the fluid nature of cell membranes, and is the first known example of a rotary motor in eukaryotic cells. The rotary motor is an axostyle that runs through the protozoan like a driveshaft and generates relative torque along its length. The entire plasma membrane of the cell is actually fluid and adheres to the underlying cytoplasm/cytoskeleton, passively following its movements.  Cell membranes clearly do not determine the shape or movements of cells. Ectosymbiotic rod bacteria attached to the cell surface are flagellated and propel the protozoan, the first known case of prokaryotic flagella driving locomotion of eukaryotic cells.}
\runningauthor{Tamm}
\jname{Journal of Trial \& Error}
\jyear{2026}
\paperdoi{10.36850/94da-4482}
\paperreceived{March 25, 2025}
\author[1]{\mbox{Sidney L. Tamm\orcid{0000-0002-4756-778X}}}
\affil[1]{Marine Biological Laboratory (MBL), Woods Hole, Massachusetts, United States}
\corremail{\href{mailto:tamm@bu.edu}{tamm@bu.edu}}
\corraddress{Boston University}
\paperaccepted{April 5, 2026}
\paperpublished{May 16, 2026}
\paperpublisheddate{2025-05-16}
\jwebsite{https://journal.trialanderror.org}



\begin{document}
\begin{frontmatter}
  \maketitle
  \begin{abstract}
    \printabstracttext
  \end{abstract}
\end{frontmatter}





	\section{Serendipitous Discovery \#1: Direct Visualization of Cell Membrane Fluidity}







	In the early 1970s, a colleague at IU who knew of my research happened to go down to Florida to visit his mother. She complained of a dead grapefruit tree in the yard that was infested with termites. My colleague was glad to saw off the tree and bring back logs full of termites to me.







	At the time I had an undergraduate in my IU lab who needed a research project to do, so I asked him to examine the hindgut protozoa from this Florida termite, \emph{Cryptotermes}\emph{ }\emph{cavifrons}. I showed him how to split the grapefruit tree logs with a hatchet to collect termites from their galleries, pull out a hindgut with forceps, put the isolated hindgut in a drop of termite Ringers solution on a microscope slide, tear it open to disperse the protozoa, and then seal the preparation with a coverslip to do microscopy of the protozoa.










	Later, when I asked the student if he was seeing anything interesting, he replied, "No, lots of flagellates swimming about and some protozoa just sitting there with one part turning around and around."







	"Let me see," I said excitedly, for I knew that single-celled protozoa are bounded by a continuous unbroken plasma membrane. Sure enough, the flagellated anterior end of a devescovinid protozoan named \emph{Caduceia}\emph{ }\emph{versatilis} was rotating continuously in the same direction (clockwise, as viewed from the anterior) relative to the rest of the cell at a speed of one rotation every 1.5 sec without stopping or reversing direction. In freshly isolated cells the body also rotated but in the opposite direction. Rotation of the head or body cytoplasm included the plasma membrane, using ectosymbiotic rod bacteria attached to the surface as visible and natural membrane markers. The continual unidirectional rotation of the membrane of the head in relation to the membrane of the body could only occur if the intervening membrane in the shear zone was extremely fluid. This accidental observation offers direct visual evidence for the fluid nature of cell membranes.







	The serendipitous discovery of this rotating head protozoan immediately became the focus of my research and I was joined by my wife and electron microscopist, Signhild Tamm. Our research efforts did not just lead to an extension of these findings, but produced further surprising and important discoveries about cells.







	Thin-section electron microscopy and freeze-fracture electron microscopy of \emph{Caduceia}\emph{ }confirmed that the trilaminar membrane structure and lipid bilayer are continuous and therefore fluid across the surface shear zone.







	Observations of additional secondary membrane shear zones during changes in cell shape indicated that the membrane between head and body is not intrinsically more fluid than the membrane of the rest of the cell. Instead, the entire membrane was shown to be potentially as fluid as the membrane between head and body. Cytochemical probes indicated no difference in sterol content between shearing and nonshearing regions of the membrane. Thus, fluidity is only expressed at certain membrane locations for geometric and/or mechanical reasons (i.e., cell shape and proximity to rotating cytoplasmic structures). The cell membrane of this protozoan may be viewed as a fluid which adheres nonspecifically to the underlying cytoplasm/cytoskeleton and passively follows its movements.







	These findings oppose the traditional model that an elastic plasma membrane determines the shape and movements of cells. I initially faced skeptical scientists, who did not believe that cells could continually rotate one part relative to another part, even when they saw it in movies or through a microscope in my lab! However, a fluid mosaic model of cell membranes with lipid and protein motility is now widely incorporated in textbooks. The discovery of a protozoan which proves this model directly, though, has not led to broader investigations of plasma membrane structure and its fluid features.







	\section{Serendipitous Discovery \#2: What Turns the Head? }







	Another novel discovery was the nature of the rotary mechanism that turns the head relative to the body of \emph{Caduceia} — the first example of a rotary motor in higher (eukaryotic) cells. Chance observations and inhibitor experiments indicated that the four flagella which arise from the head do not turn it. Instead, the motor is internal and consists of a cytoskeletal rod (axostyle) which runs through the cell like a drive-shaft from the head to the posterior end, and uses ATP as the energy source. Laser microbeam dissection of the axostyle at different locations showed that the axostyle generates constant torque along its length. Surprisingly, the molecular basis of force production by this rotary motor has still not been determined. Another question which has not been followed up is: Why do parts of the cell rotate at all? What is the function of opposing rotations of the head and body?
	







	But we do know that the rotational motility of \emph{Caduceia}\emph{ }\emph{versatilis} is not unique. Other related devescovinid protozoans in Australia show similar rotational motility where the head rotates relative to the body, providing additional visual evidence for the fluid nature of cell membranes. This was first reported by the protozoologist Harold Kirby in the 1940s, so I went to Australia and confirmed his findings. Moreover, I observed that the direction of head rotation in these devescovinids from the Southern Hemisphere was also clockwise as viewed anteriorly.



	\section{Serendipitous Discovery \#3: Bacterial Flagella Propel \emph{Caduceia} }







	In densely packed protozoa from freshly opened hindguts, I had routinely noticed (and ignored) vigorous slithering movements of devescovinids through the crowded mass. When devescovinids emerged from the mass and were no longer surrounded by other protozoa, they slowed down and stopped gliding, even though rotational movements of the head and its beating flagella continued actively. Locomotory movements reappeared if these stationary cells were gently flattened by the coverslip so most of their body surface was in contact with a substrate.







	I had also noticed that the 2,000 - 3,000 ectosymbiotic rod bacteria attached to the head and body were aligned end-to-end in parallel rows of specialized open pockets of the cell membrane. Surprisingly, the rod bacteria had bacterial flagella on their exposed side, even though they normally never leave the surface of the protozoan host. Flagella of bacteria (called prokaryotes), which do not have a true nucleus, are much simpler than microtubular flagella of eukaryotic cells with nuclei. Bacterial flagella are bare helical filaments not enclosed by an extension of the cell membrane, that can rotate in opposite directions by a tiny rotary motor and form propulsive helical bundles, using a proton gradient (not ATP) as energy.







	It then dawned on me that the flagella of the attached rod bacteria must be involved in the vigorous contact-induced gliding locomotion of \emph{Caduceia}. So I performed extensive experiments to test this idea. I found that (1) gliding continues when the devescovinid's own flagella and rotary axostyle are inactivated; (2) agents which inhibit bacterial flagellar motility, but not the protozoan's motile systems, stop gliding movements; (3) isolated vesicles derived from the surface of the devescovinid rotate at speeds dependent on the number of rod bacteria still attached; (4) individual rod bacteria can move independently over the surface of compressed cells; and (5) wave propagation by the flagellar bundles of the bacteria can be seen by video-enhanced polarization microscopy.







	Proximity to solid boundaries (other cells or surfaces) may be required to align the flagellar bundles of adjacent bacteria in the same direction, and/or increase their propulsive efficiency (wall effect).







	Locomotion of \emph{Caduceia} is thus powered not by the cell's own flagella, nor by its rotary axostyle, but by the flagella of thousands of rod bacteria which live on its surface. Since bacterial flagella rotate, an additional novelty of this system is that the surface bearing the prokaryotic rotary motors is turned by the eukaryotic rotary motor within. These ectosymbiotic rod bacteria on \emph{Caduceia} have recently been identified and named (not by me) as a new genus and species, \emph{Tammella}\emph{ }\emph{caduceiae}\emph{. }







	\section{Conclusion}







	I accidentally found this remarkable protozoan that I certainly was not looking for. A series of fortunate chance encounters and situations, with occasional flashes of sagacity, led to three serendipitous discoveries about protozoa and cells in general: (1) cell membranes are visually fluid and allow unidirectional rotational shear between different regions; (2) the rotary motor in the best studied protozoan is a cytoskeletal drive-shaft that generates constant relative torque along its length; and (3) flagella of thousands of ectosymbiotic rod bacteria can propel the host cell. These unexpected discoveries have changed our understanding of how prokaryotic and eukaryotic cells function and can live beneficially together. Readers may be interested in finding more information about the phenomena discussed here in the citations of my published work in Further Reading.











	\section{Further Reading}



	



	d'Ambrosio, U., Dolan, M., Wier, A. M., \& Margulis, L. (1999). Devescovinid trichomonad with axostyle-based rotary motor (“Rubberneckia”): Taxonomic assignment as \emph{Caduceia}\emph{ }\emph{versatilis} sp. nov.. \emph{European Journal of Protistology, 35}(3), 327--337. \url{https://doi.org/10.1016/s0932-4739(99)80011-x}



	



	Hongoh, Y., Sato, T., Dolan, M. F., Noda, S., Ui, S., Kudo, T., \& Ohkuma, M. (2007). The motility symbiont of the termite gut flagellate \emph{Caduceia}\emph{ }\emph{versatilis} is a member of the “synergistes” group. \emph{Applied and Environmental Microbiology, 73}(19), 6270--6276. \url{https://doi.org/10.1128/aem.00750-07}



	



	Tamm, S. L. (1978). Laser microbeam study of a rotary motor in termite flagellates.	Evidence that the axostyle complex generates torque. \emph{The Journal of Cell}	\emph{Biology, 78}(1), 76--92. \url{https://doi.org/10.1083/jcb.78.1.76}



	



	Tamm, S. L. (1979). Membrane movements and fluidity during rotational motility of a termite flagellate. A freeze-fracture study. \emph{Journal of Cell Biology, 80}(1), 141--149. \url{https://doi.org/10.1083/jcb.80.1.141}



	



	Tamm, S. L. (1980). The ultrastructure of prokaryotic-eukaryotic cell junctions. \emph{Journal of Cell Science, 44}(1), 335--352. \url{https://doi.org/10.1242/jcs.44.1.335}


	



	Tamm, S. L. (1982). Flagellated ectosymbiotic bacteria propel a eucaryotic cell. \emph{Journal} \emph{of Cell Biology, 94}(3), 697--709. \url{https://doi.org/10.1083/jcb.94.3.697}



	



	Tamm, S. L. (2008). Unsolved motility looking for answer. \emph{Cell Motility, 65}(6), 435--440. \url{https://doi.org/10.1002/cm.20276}



	



	Tamm, S. L., \& Tamm, S. (1974). Direct evidence for fluid membranes. \emph{Proceedings of} \emph{the National Academy of Sciences, 71}(11), 4589--4593. \url{https://doi.org/10.1073/pnas.71.11.4589}



	



	Tamm, S. L., \& Tamm, S. (1976). Rotary movements and fluid membranes in termite flagellates. \emph{Journal of Cell Science, 20}(3), 619--638. \url{https://doi.org/10.1242/jcs.20.3.619}



	



	Tamm, S. L., \& Tamm, S. (1983). Distribution of sterol-specific complexes in a	continually shearing region of a plasma membrane and at procaryotic-eucaryotic cell junctions. \emph{The Journal of Cell Biology, 97}(4), 1098--1106. \url{https://doi.org/10.1083/jcb.97.4.1098}



	



	Yamin, M. A., \& Tamm, S. L. (1982). ATP reactivation of the rotary axostyle in termite	flagellates: effects of dynein ATPase inhibitors. \emph{The Journal of Cell Biology, 95}(2), 589--597. \url{https://doi.org/10.1083/jcb.95.2.589}






\end{document}